\section{When can a maintainer accept a repair?}

\label{sec:introduction}

\subsection{The caller who still gets the wrong answer}

A maintainer reviewing a repair to Ekubo could ask the library for a price history and get the expected answer. A caller asking for that history through the direct interface could still get an incorrect answer. Both observations concerned the same audited finding. The library result gave the maintainer evidence about a real use of the software. Whether it warranted accepting the patch depended on what the caller was entitled to receive. The development records in Table~\ref{tab:development} show what happened when the assessment included another way to observe the reported defect.

Consider what a passing regression test would tell the maintainer in this situation. If the test checks the value returned to the library, its verdict concerns the observed value on that call. If the test instead detects the old violation, the disappearance of that signal requires interpretation. Did the operation run and give a permitted result? Could the patch have prevented it from running? Even a convincing answer about this operation leaves a further question about a caller with another documented requirement. Each question calls for evidence about a particular behaviour. The distinction matters both to the person accepting an individual patch and to anyone comparing repair systems by the fraction of patches that pass.

Checking what ordinary users can still do has a precedent in Ethereum patching. EVMPatch evaluated changes against intended functionality using recorded transactions, and Bobadilla et al. examined whether automated smart-contract fixes preserve functionality. EVMbench also combines project tests with checks withheld from the agent. These studies motivate looking beyond the example in a repair ticket. A maintainer must still decide what a particular combination of observations establishes about the submitted patch.~\cite{rodler2021evmpatch,bobadilla2026fixes,wang2026evmbench}

\subsection{Three questions about the same patch}

Research question 1 (RQ1) asks how often an assigned repair yields a patch that builds, restores the reported behaviour, and preserves a specified ordinary use of the affected feature. Research question 2 (RQ2) asks how often those same attempts also restore a distinct, withheld behaviour, among findings with sufficient evidence for that additional requirement. Research question 3 (RQ3) asks whether changing which example an agent sees changes which behaviour it repairs. An agent might usually repair the library behaviour under either assignment. To establish the effect of showing it, we would need to assess that behaviour even on patches produced when the agent saw a different example.

We contribute an analysis of archived attempts on two findings, an account of which audited findings have enough evidence to grade, and a registered design for comparing two assignments within the same finding. The first contribution concerns patches that agents actually produced. The second identifies where audit reports, historical revisions and their tests suffice to specify a repair, and where they leave a maintainer without a defensible grade. By 26 November one case had reached human review for admission. No case was signed and no repair session in the expanded experiment had started (Table~\ref{tab:qualification}). The third contribution therefore defines the observations needed to answer the three questions. A maintainer could use each answer differently. The first concerns acceptance of the ticket, the second the coverage of that decision, and the third how to write the ticket itself.
